\documentclass[7pt]{article}
\usepackage{graphicx}
\usepackage{xcolor}
\usepackage{float}
\usepackage[margin=1in]{geometry}
\usepackage{url} %COLOCAR LINKS - insertar este paquete
\usepackage{hyperref}
\usepackage[utf8]{inputenc}

\title{\textbf{Algoritmos y Sistemas Computacionales: \\ Nacho te recomienda.}}
\author{Ignacio Pizarro}
\date{\textit{\underline{12 de Mayo, 2025.}}}


\begin{document}

\maketitle

\begin{figure}[H]
    \centering
    \includegraphics[width=0.8\textwidth]{Nintendo-DS-Lite-Black-Open.png}
    \caption{"Una imagen molona de la Nintendo DS Lite."}
\end{figure}

\newpage


\Large\section{\textbf{¿Por qué existe?}}
\hfill
\hfill
\\

Este programa .py existe única y exclusivamente por el aburrimiento del autor a la hora de iniciar su computadora y abrir YouTube, y lo primero que salio en su pagina principal era un video de "Jugando un mes unicamente juegos de Nintendo DS", y se le ocurrió hacer una ruleta en Python, la cual sorteara entre los juegos que se irían añadiendo a una lista. (No fue así).
\\
\\
Al menos la idea de la Nintendo DS se mantuvo, y el programa funciona tal que así:
\\
\\
A la hora de iniciar el programa, nos pedirá un Input (El cual debe ser el nombre de una consola <Que hasta el momento es unicamente la Nintendo DS>). En el momento que introducimos el nombre de la consola, el programa nos dará una pequeña lista de juegos de dicha consola. El siguiente paso sera anotar el nombre del juego tal y como sale en la lista, o si no el programa no funcionara.\\
\\
Luego de esto, si se cumplieron todos los pasos de manera correcta, el programa soltara lo que piensa el autor de este archivo LaTeX sobre el juego. (Algunos juegos aun no están  programados y se irán agregando con las versiones posteriores.)
\newline
\newline


\begin{figure}
    \centering
    \includegraphics[width=0.1\linewidth]{images.png}
    \caption{:D}
    \label{fig:enter-label}
\end{figure}

\end{document}
